\documentclass[letterpaper,10pt]{article}

\usepackage{latexsym}
\usepackage[empty]{fullpage}
\usepackage{titlesec}
\usepackage{marvosym}
\usepackage[usenames,dvipsnames]{color}
\usepackage{verbatim}
\usepackage{enumitem}
\usepackage[hidelinks]{hyperref}
\usepackage{fancyhdr}
\usepackage[brazilian]{babel}
\usepackage{tabularx}
\input{glyphtounicode}

\usepackage[sfdefault]{roboto}

\pagestyle{fancy}
\fancyhf{}
\fancyfoot{}
\renewcommand{\headrulewidth}{0pt}
\renewcommand{\footrulewidth}{0pt}
\setlength{\footskip}{5pt}

\addtolength{\oddsidemargin}{-0.5in}
\addtolength{\evensidemargin}{-0.5in}
\addtolength{\textwidth}{1in}
\addtolength{\topmargin}{-.5in}
\addtolength{\textheight}{1.0in}

\urlstyle{same}
\raggedbottom
\setlength{\tabcolsep}{0in}

% Configuração para justificar o texto em blocos e sem recuo
\setlength{\parindent}{0pt}
\setlength{\parskip}{10pt}

\titleformat{\section}{\Large\bfseries\scshape\raggedright}{}{0em}{}[\titlerule]

\pdfgentounicode=1

\newcommand{\resumeItem}[1]{\item\small{#1}}
\newcommand{\resumeSubheading}[4]{
\vspace{-1pt}\item
  \begin{tabular*}{0.97\textwidth}[t]{l@{\extracolsep{\fill}}r}
    \textbf{#1} & #2 \\
    \textit{#3} & \textit{#4} \\
  \end{tabular*}\vspace{-7pt}
}
\renewcommand\labelitemii{$\vcenter{\hbox{\tiny$\bullet$}}$}
\newcommand{\resumeSubHeadingList}{\begin{itemize}[leftmargin=0.15in, label={}]}
\newcommand{\resumeSubHeadingListEnd}{\end{itemize}}

% === PERSONALIZAR: Altere o cargo e a empresa para cada candidatura ===
\newcommand{\targetRole}{Estágio em Desenvolvimento}
\newcommand{\targetCompany}{[Nome da Empresa]}

\begin{document}

\begin{center}
  \textbf{\Huge Gabriel Frigo} \\ \vspace{5pt}
  \small +55 11 95570-2000 $|$ \href{mailto:gabriel.frigo4@gmail.com}{gabriel.frigo4@gmail.com} $|$
  \href{https://linkedin.com/in/gabriel-frigo-b6727b275/}{linkedin.com/in/gabriel-frigo} $|$
  \href{https://github.com/GabrielFrigo4}{github.com/GabrielFrigo4} $|$
  \href{https://gabrielfrigo.dev.br}{gabrielfrigo.dev.br}
\end{center}

\vspace{10pt}

\textbf{Assunto: Candidatura à vaga de \targetRole}

Prezada Equipe de Recrutamento da \targetCompany,

Escrevo para expressar meu interesse na vaga de \textbf{\targetRole}. Sou estudante de \textbf{Ciência da Computação na UFABC}, com um perfil que une programação de sistemas de baixo nível a uma visão ampla de arquitetura de software, administração de servidores e infraestrutura.

No lado prático, tenho experiência com \textbf{APIs RESTful em Python 3} para serviços financeiros (QI Tech), \textbf{deploy e administração de servidores Linux} na Oracle Cloud, e conteinerização com \textbf{Docker, Podman e LXC}. Desenvolvi um \textbf{servidor HTTP concorrente em C} do zero — compatível com GNU/Linux e FreeBSD — e atualmente mantenho um catálogo de lojas em produção com \textbf{Go e PocketBase}. Além disso, domino linguagens de baixo nível como \textbf{C, C++, Rust e Assembly}, que aplico na administração de projetos no \textbf{Green Team Hacker Club} e no desenvolvimento de software para robótica competitiva.

Minha base analítica é reforçada por pesquisas de Iniciação Científica em \textbf{Fluxos em Redes (CNPq)} e \textbf{Teoria de Ramsey (PICME)}, além de resultados em competições como a \textbf{Medalha de Ouro e 13ª posição nacional na OBMEP} e a participação na \textbf{Fase Regional da Maratona SBC de Programação (ICPC)}. Sou entusiasta de \textbf{padrões de projeto, programação funcional e otimização algorítmica} — e atuo como mentor de algoritmos para a OBI.

Acredito que minha versatilidade técnica — transitando entre infraestrutura, backend e sistemas de alto desempenho — me qualifica para contribuir de forma efetiva com os times da \targetCompany. Estou à disposição para uma entrevista.

Atenciosamente,

\textbf{Gabriel Frigo}

\end{document}
